\subsection[Preliminary definitions for p=3]{Preliminary definitions for $\boldsymbol{p=3}$}

Next, we define the indecomposable and irreducible representations of $S_3$.

\begin{Proposition}[\cite{Alperin_1986}]
 There are $6$ indecomposable or irreducible representations in $S_3$ in characteristic $3$:
 \begin{enumerate}\itemsep=0pt
 \item[$1$.] $\triv$ is the irreducible representation of $S_3$ that is acted on trivially by $S_3$.
 \item[$2$.] $\sign$ is the irreducible representation of $S_3$ that is acted on by negation by the transpositions.
 \item[$3$.] $\siv$ is the indecomposable representation that contains a copy of $\triv$ as a subrepresentation, such that the quotient of $\siv$ by this subrepresentation is $\sign$.
 \item[$4$.] $\tign$ is the indecomposable representation that contains a copy of $\sign$ as a subrepresentation, such that the quotient of $\tign$ by this subrepresentation is $\triv$.
 \item[$5$.] $\trix$ is the indecomposable representation that contains a copy of $\siv$ as a subrepresentation, such that the quotient of $\trix$ by this subrepresentation is $\triv$.
 \item[$6$.] $\six$ is the indecomposable representation that contains a copy of $\tign$ as a subrepresentation, such that the quotient of $\six$ by this subrepresentation is $\sign$.
 \end{enumerate}
\end{Proposition}

Provided are some examples of copies of these indecomposable representations:

\begin{Example}\label{exSiv}
 The space $W\subset\F_3[x_1,x_2,x_3]$ spanned by $x_1+x_2+x_3$ and $x_1-x_2$ over $\mathbf{F}_3$ is copy of $\siv$. Indeed, the space $T\subset W$ spanned by $x_1+x_2+x_3$ is a copy of $\triv$. One can check $x_1-x_2\in W/T$ is acted by negation by all transpositions in $S_3$ and $W/T$ is $1$-dimensional so $W/T$ is a copy of $\sign$. Finally, it is easy to show that there are no copies of $\triv$ or $\sign$ in $W$ other than $T$. Since $V$ has a unique irreducible subrepresentation, it is indecomposable, and we conclude that it is a copy of $\siv$.
\end{Example}

\begin{Example}\label{deg1polys}
 The space $V\subset\F_3[x_1,x_2,x_3]$ consisting of homogeneous linear polynomials is a~copy of $\trix$. Indeed, $W\subset V$ from Example~\ref{exSiv} is a copy of $\siv$. Then~$V/W$ is one-dimensional, and one can check that it is a copy of $\triv$. Finally, it is easy to show that there are no copies of $\triv$ or $\sign$ in $V$ other than $T$, so $V$ has a unique irreducible subrepresentation, it is indecomposable, and we conclude that it is a copy of $\trix$.
\end{Example}

\begin{Example}
 Similarly, one may check that the space $U$ spanned by
 \[(x_1-x_2)(x_1-x_3)(x_2-x_3),\qquad-x_1^2x_2-x_1^2x_3+x_1x_2^2+x_1x_3^2\]
 over $\mathbf{F}_3$ is a copy of $\tign$ and that the space spanned by
 \[(x_1-x_2)(x_1-x_3)(x_2-x_3),\qquad-x_1^2x_2-x_1^2x_3+x_1x_2^2+x_1x_3^2,\qquad(x_1-x_2)x_1x_2\]
 is a copy of $\six$.
\end{Example}

Similarly to the $p=2$ case, we define $Q_m(3,\mathbf{F}_3)_{\sign}$ and $Q_m(3,\mathbf{F}_3)_{\triv}$ to be the direct sum of all copies of $\sign$ and $\triv$ in $Q_m(3,\mathbf{F}_3)$, respectively.
\begin{Proposition}[\cite{wang2022explicit}]\label{prop:trivsign}
 As $\mathbf{F}_3[x_1,x_2,x_3]^{S_3}$-modules,
 \begin{enumerate}\itemsep=0pt
 \item[$1$.] $Q_m(3,\mathbf{F}_3)_\triv$ is freely generated by $1$.
 \item[$2$.] $Q_m(3,\mathbf{F}_3)_\sgn$ is freely generated by $\prod_{i_1<i_2}(x_{i_1}-x_{i_2})^{2m+1}$.
 \end{enumerate}
\end{Proposition}

Next we define $Q_m(3,\mathbf{F}_3)_{\siv}$ as the direct sum of all copies of $\sign$, $\triv$, and $\siv$. For this paper we consider generators of $Q_m(3,\mathbf{F}_3)_{\siv}$ to be homogeneous polynomials other than $1$ and \smash{$\prod_{i_1<i_2}(x_{i_1}-x_{i_2})^{2m+1}$} such that they are in the $(-1)$-eigenspace of $s_{12}$ and are in a~generators and relations presentation of $Q_m(3,\mathbf{F}_3)_{\siv}$ as an $\mathbf{F}_3[x_1,x_2,x_3]^{S_3}$-module with the least number of generators. Moreover, if $K$ is a generator of $Q_m(3,\mathbf{F}_3)_{\siv}$ then it necessarily generates a copy of $\siv$ since we assumed $K$ neither generates $\triv$ nor $\sign$.

\begin{Remark}
 Similar to in the $p=2$ case, we cannot define $Q_m(3,\mathbf{F}_3)_{\siv}$ to exclude copies of $\sign$ since we can add elements of $Q_m(3,\mathbf{F}_3)_{\sign}$ to copies of $\siv$ and still obtain a copy of $\siv$. For example, the spaces spanned by
 \[\big(x_1^6-x_2^6\big)(x_1+x_2+x_3)^3,\qquad\big(x_1^6+x_2^6+x_3^6\big)(x_1+x_2+x_3)^3\]
 and
 \[\prod_{i_1<i_2}(x_{i_1}-x_{i_2})^{3}+\big(x_1^6-x_2^6\big)(x_1+x_2+x_3)^3,\qquad\big(x_1^6+x_2^6+x_3^6\big)(x_1+x_2+x_3)^3\]
 generate two copies of $\siv$ in $Q_1(3,\F_3)$, and their sum contains
 \[\prod_{i_1<i_2}(x_{i_1}-x_{i_2})^{3}\in Q_1(3,\mathbf{F}_3)_{\sign}.\]
\end{Remark}

\begin{Remark}
 One could define subspaces of $Q_m(3,\mathbf{F}_3)$ for $\trix$, $\six$, $\tign$ similar to $Q_m(3,\mathbf{F}_3)_{\siv}$, however this is not particularly helpful, as unlike for ${p=2}$, we cannot decompose $Q_m(3,\mathbf{F}_3)$ into a direct sum of subspaces of this form. The space $Q_m(3,\mathbf{F}_3)_{\siv}$ is still relevant, as it is the critical piece to understanding quasi-invariants in characteristic 3, as we see in Sections \ref{sec:wang} and \ref{sec:renxu}.
\end{Remark}
